% Options for packages loaded elsewhere
\PassOptionsToPackage{unicode}{hyperref}
\PassOptionsToPackage{hyphens}{url}
\documentclass[
  11pt,
]{article}
\usepackage{xcolor}
\usepackage[margin=1in]{geometry}
\usepackage{amsmath,amssymb}
\setcounter{secnumdepth}{-\maxdimen} % remove section numbering
\usepackage{iftex}
\ifPDFTeX
  \usepackage[T1]{fontenc}
  \usepackage[utf8]{inputenc}
  \usepackage{textcomp} % provide euro and other symbols
\else % if luatex or xetex
  \usepackage{unicode-math} % this also loads fontspec
  \defaultfontfeatures{Scale=MatchLowercase}
  \defaultfontfeatures[\rmfamily]{Ligatures=TeX,Scale=1}
\fi
\usepackage{lmodern}
\ifPDFTeX\else
  % xetex/luatex font selection
  \setmainfont[BoldFont=NotoSerif-Bold.ttf,ItalicFont=NotoSerif-Italic.ttf,BoldItalicFont=NotoSerif-BoldItalic.ttf]{NotoSerif-Regular.ttf}
  \setmathfont[]{Latin Modern Math}
\fi
% Use upquote if available, for straight quotes in verbatim environments
\IfFileExists{upquote.sty}{\usepackage{upquote}}{}
\IfFileExists{microtype.sty}{% use microtype if available
  \usepackage[]{microtype}
  \UseMicrotypeSet[protrusion]{basicmath} % disable protrusion for tt fonts
}{}
\makeatletter
\@ifundefined{KOMAClassName}{% if non-KOMA class
  \IfFileExists{parskip.sty}{%
    \usepackage{parskip}
  }{% else
    \setlength{\parindent}{0pt}
    \setlength{\parskip}{6pt plus 2pt minus 1pt}}
}{% if KOMA class
  \KOMAoptions{parskip=half}}
\makeatother
\usepackage{longtable,booktabs,array}
\newcounter{none} % for unnumbered tables
\usepackage{calc} % for calculating minipage widths
% Correct order of tables after \paragraph or \subparagraph
\usepackage{etoolbox}
\makeatletter
\patchcmd\longtable{\par}{\if@noskipsec\mbox{}\fi\par}{}{}
\makeatother
% Allow footnotes in longtable head/foot
\IfFileExists{footnotehyper.sty}{\usepackage{footnotehyper}}{\usepackage{footnote}}
\makesavenoteenv{longtable}
\setlength{\emergencystretch}{3em} % prevent overfull lines
\providecommand{\tightlist}{%
  \setlength{\itemsep}{0pt}\setlength{\parskip}{0pt}}
\setcounter{secnumdepth}{-\maxdimen}
\usepackage{longtable,booktabs,array}
\setlength{\emergencystretch}{6em}
\tolerance=1500
\usepackage{bookmark}
\IfFileExists{xurl.sty}{\usepackage{xurl}}{} % add URL line breaks if available
\urlstyle{same}
\hypersetup{
  hidelinks,
  pdfcreator={LaTeX via pandoc}}

\author{}
\date{}

\begin{document}

\section{Reading the Indus Seals: A Translation Corpus of 7,138
Inscriptions}\label{reading-the-indus-seals-a-translation-corpus-of-7138-inscriptions}

\textbf{Tristen Kyle Pierson}\\
Glossa-Lab / BitConcepts LLC\\
May 2026\\
Companion to: DOI 10.5281/zenodo.20381178

\begin{center}\rule{0.5\linewidth}{0.5pt}\end{center}

\subsection{Overview}\label{overview}

This report presents the first mass translation of Indus Valley seal
inscriptions using the 605-sign Proto-Dravidian anchor table (Pierson
2026). All readings are computational proposals, not confirmed
translations, and should be treated as hypotheses pending peer review.

{\def\LTcaptype{none} % do not increment counter
\begin{longtable}[]{@{}ll@{}}
\toprule\noalign{}
Metric & Value \\
\midrule\noalign{}
\endhead
\bottomrule\noalign{}
\endlastfoot
Total inscriptions translated & 7,138 \\
Corpora & Holdat (1,669 seals) + Yajnadevam (5,469 inscriptions) \\
Mean sign coverage & 87.9\% \\
Fully translated (100\% coverage) & 5,372 (75.3\%) \\
Archaeological sites represented & 77 \\
Unique inscription formulas & 3,352 \\
PDr vs Sanskrit side-by-side & 3,731 comparisons \\
\end{longtable}
}

\begin{center}\rule{0.5\linewidth}{0.5pt}\end{center}

\subsection{The Inscriptions: What Do They
Say?}\label{the-inscriptions-what-do-they-say}

The most common seal readings confirm the guild-identity model proposed
in the preprint. Seals encode short formulaic titles --- not commodity
tallies, not religious texts, but professional identity markers
analogous to medieval guild marks.

\subsubsection{Most Common Seal
Formulas}\label{most-common-seal-formulas}

{\def\LTcaptype{none} % do not increment counter
\begin{longtable}[]{@{}
  >{\raggedright\arraybackslash}p{(\linewidth - 4\tabcolsep) * \real{0.3333}}
  >{\raggedright\arraybackslash}p{(\linewidth - 4\tabcolsep) * \real{0.3333}}
  >{\raggedright\arraybackslash}p{(\linewidth - 4\tabcolsep) * \real{0.3333}}@{}}
\toprule\noalign{}
\begin{minipage}[b]{\linewidth}\raggedright
Reading
\end{minipage} & \begin{minipage}[b]{\linewidth}\raggedright
Count
\end{minipage} & \begin{minipage}[b]{\linewidth}\raggedright
Interpretation
\end{minipage} \\
\midrule\noalign{}
\endhead
\bottomrule\noalign{}
\endlastfoot
kur & 284 (4.0\%) & ``Hill/short'' --- single-sign administrative
marker \\
min/min & 150 (2.1\%) & ``Fish/star'' --- fish guild or astronomical
marker \\
kur-kur & 136 (1.9\%) & Doubled hill marker --- emphatic or plural \\
kur-puli & 130 (1.8\%) & ``Hill-tiger'' --- tiger-clan hill guild \\
kur-tol & 112 (1.6\%) & ``Hill-shoulder'' --- craft guild title \\
kur-koL & 96 (1.3\%) & ``Hill-vessel'' --- jar/pottery guild \\
koL & 60 (0.8\%) & ``Vessel/jar'' --- single-sign potter marker \\
puli & 41 (0.6\%) & ``Tiger'' --- tiger clan \\
nallavar & 40 (0.6\%) & ``Good people'' --- elite/noble title \\
al-kur-kat & 39 (0.5\%) & ``Man-hill-bond'' --- compound administrative
title \\
\end{longtable}
}

The dominant pattern is
\textbf{{[}ANIMAL/CRAFT{]}-{[}PLACE/TITLE{]}-{[}SUFFIX{]}} --- exactly
the tripartite INITIAL-MEDIAL-TERMINAL grammar confirmed at 6.3x above
null across 76 sites.

\begin{center}\rule{0.5\linewidth}{0.5pt}\end{center}

\subsection{The Indus Valley Civilization: What the Seals
Reveal}\label{the-indus-valley-civilization-what-the-seals-reveal}

\subsubsection{A Civilization of Guilds}\label{a-civilization-of-guilds}

The translated inscriptions paint a picture of a highly organized
guild-based society. The IVC (ca. 2600-1900 BCE) was not merely a
collection of farming settlements --- it was an urban civilization with
professional specialization, standardized administrative practices, and
a continent-spanning trade network.

\textbf{Key cultural insights from the translation corpus:}

\begin{enumerate}
\def\labelenumi{\arabic{enumi}.}
\item
  \textbf{Professional identity seals.} The overwhelming majority of
  inscriptions encode guild membership and personal identity, not
  commodity counts or religious dedications. A typical seal reads: ``the
  tiger-clan hill craftsman'' or ``the vessel-guild man.'' This is
  directly analogous to later South Asian guild-seal traditions
  documented in the Arthasastra (4th century BCE).
\item
  \textbf{Animal-clan system.} The most frequent INITIAL signs are
  animal names: puli (tiger), erutu (bull), kur (hill-creature), min
  (fish). Each animal appears to designate a professional clan or guild
  --- the ``tiger clan,'' the ``bull guild,'' the ``fish-star guild.''
  This totem-based clan system is consistent with Dravidian kinship
  structures attested in historical Tamil Sangam literature.
\item
  \textbf{Standardized across 77 sites.} The same formulaic patterns
  appear at Harappa, Mohenjo-daro, Dholavira, Lothal, and 73 other sites
  spanning from the Afghan border to Gujarat. This uniformity implies a
  centralized administrative system with shared scribal conventions ---
  remarkable for a civilization covering 1.5 million square kilometers.
\end{enumerate}

\subsubsection{Sites and Their
Inscriptions}\label{sites-and-their-inscriptions}

{\def\LTcaptype{none} % do not increment counter
\begin{longtable}[]{@{}lll@{}}
\toprule\noalign{}
Site & Inscriptions & Region \\
\midrule\noalign{}
\endhead
\bottomrule\noalign{}
\endlastfoot
Harappa & 2,582 & Punjab (Pakistan) \\
Mohenjo-daro & 1,896 & Sindh (Pakistan) \\
Dholavira & 230 & Gujarat (India) \\
Lothal & 200 & Gujarat (India) \\
Kalibangan & 195 & Rajasthan (India) \\
Chanhu-daro & 73 & Sindh (Pakistan) \\
Nausharo & 38 & Balochistan (Pakistan) \\
Banawali & 26 & Haryana (India) \\
Lakhanjo-daro & 14 & Sindh (Pakistan) \\
Rakhigarhi & 13 & Haryana (India) \\
\end{longtable}
}

The two great cities --- Harappa and Mohenjo-daro --- account for 63\%
of all inscriptions. Dholavira, the third-largest site, is notable for
its monumental signboard (the only known large-scale Indus inscription)
and its distinctive seal repertoire.

\subsubsection{Materials: The Technology of
Seals}\label{materials-the-technology-of-seals}

{\def\LTcaptype{none} % do not increment counter
\begin{longtable}[]{@{}lll@{}}
\toprule\noalign{}
Material & Count & Notes \\
\midrule\noalign{}
\endhead
\bottomrule\noalign{}
\endlastfoot
Steatite & 2,549 & Primary seal material --- soft stone, easy to
carve \\
Clay & 913 & Sealings (impressions) and tablets \\
Faience & 596 & Glazed ceramic --- luxury goods markers \\
Copper & 261 & Copper tablets --- administrative records \\
Terracotta & 76 & Fired clay --- permanent records \\
Ivory & 36 & Rare luxury --- high-status seals \\
\end{longtable}
}

Steatite dominates (46\% of all inscribed objects), consistent with its
use as the standard seal material throughout the IVC. The presence of
261 copper tablets is significant --- copper was expensive, suggesting
these were official administrative records rather than personal seals.

\subsubsection{Chronological
Distribution}\label{chronological-distribution}

The corpus spans the full IVC timeline:

\begin{itemize}
\tightlist
\item
  \textbf{Early period}: 49 inscriptions --- the script's formative
  phase
\item
  \textbf{Intermediate/Period 3}: 1,392 inscriptions --- the mature
  Harappan peak
\item
  \textbf{Late period}: 327 inscriptions --- the decline phase before
  \textasciitilde1900 BCE
\end{itemize}

The same formulaic patterns persist across all periods, suggesting
remarkable scribal conservatism over 700+ years.

\begin{center}\rule{0.5\linewidth}{0.5pt}\end{center}

\subsection{Proto-Dravidian vs Sanskrit:
Side-by-Side}\label{proto-dravidian-vs-sanskrit-side-by-side}

The translation corpus includes 3,731 inscriptions where both our
Proto-Dravidian reading and Yajnadevam's Sanskrit reading are available.
As reported in the preprint, agreement is 0/34 for directly comparable
signs --- the two hypotheses produce completely different readings for
the same inscriptions.

This is not merely a disagreement about phonetic values --- the entire
semantic framework differs. Our readings produce guild titles and
professional identifiers; the Sanskrit hypothesis produces religious
invocations and deity names. The archaeological context (commercial
seals, administrative tablets, port-city trade goods) strongly favors
the guild-identity interpretation.

\begin{center}\rule{0.5\linewidth}{0.5pt}\end{center}

\subsection{Limitations}\label{limitations}

\begin{enumerate}
\def\labelenumi{\arabic{enumi}.}
\tightlist
\item
  \textbf{Coverage gap.} 12.1\% of sign tokens remain untranslated (66
  Yajnadevam glyph IDs with no Mahadevan equivalent).
\item
  \textbf{No confirmed translations.} All readings are computational
  proposals. The 605-sign anchor table has not been independently
  validated by epigraphists.
\item
  \textbf{Single-hypothesis.} All readings assume Proto-Dravidian. A
  formal comparison against Proto-Munda readings would strengthen or
  weaken this interpretation.
\item
  \textbf{No bilingual text.} Until a bilingual inscription is found,
  all translations remain probabilistic.
\end{enumerate}

\begin{center}\rule{0.5\linewidth}{0.5pt}\end{center}

\subsection{Data Availability}\label{data-availability}

All translation data is open source:

\begin{itemize}
\tightlist
\item
  Full corpus: \texttt{outputs/seal\_translations.json} and
  \texttt{.csv}
\item
  Cultural summary: \texttt{outputs/seal\_translations\_summary.json}
\item
  PDr vs Sanskrit: \texttt{outputs/pdr\_vs\_sanskrit.csv}
\item
  Anchor table: \texttt{backend/reports/INDUS\_FINAL\_ANCHORS.json}
\item
  Repository: https://github.com/BitConcepts/glossa-lab
\item
  Preprint DOI: https://zenodo.org/records/20381178
\end{itemize}

\begin{center}\rule{0.5\linewidth}{0.5pt}\end{center}

\emph{Glossa-Lab, May 2026}\\
\emph{Tristen Kyle Pierson}

\end{document}
